\section{How to Create a Testing Project}\label{sec:project}

In \toolname the testing activity requires the creation of a
testing project. A testing project is characterized by a file
storing the necessary information about (i) the base class file,
(ii) the complete set of classes required by the base class, (iii)
the set of classes to be instrumented (tested), and (iv) the set
of classes that are not under testing. Additional information,
such as the \pk{CLASSPATH} environment variable necessary to run
the base class is also stored in the project file, whose extension
is \prjext. During the execution of any \pk{.class} file that
belongs to the set of classes under testing of a given project,
dynamic control-flow information (execution trace) is collected
and saved in a separate file that has the same name of the project
file but with a different extension (\trcext).

For example, considering the vending machine example, described in
Section~\ref{sec:example}, the \pk{vending} package is composed by
three \bci{.java} files: \pk{VendingMachine.java},
\pk{Dispenser.java} and \pk{TestDriver.java}. Suppose that these
files are saved in a directory named \verb+~\example+. The
directory structure is like:

\begin{cmd}
auri@AURIMRV ~
$ ls -l example/*
-rw-r--r--    1 auri         1313 Aug  6 09:07 Dispenser.java
-rw-r--r--    1 auri         1340 Aug  6 09:07 TestDriver.java
-rw-r--r--    1 auri          923 Aug  6 09:07 VendingMachine.java
\end{cmd}

To compile such an application one of the following command can be
used:

\begin{cmd}
auri@AURIMRV ~
$ javac -d example example/*.java
\end{cmd}

or

\begin{cmd}
auri@AURIMRV ~
$ javac -g -d example example/*.java
\end{cmd}

Observe that in the later command, the debug option is activated,
thus, the generated \pk{.class} files contains more information,
such as the real variable names. In this report we are using class
files compiled with the debug option.

After the Java source files have been compiled, the
\verb+~\example+ directory contains the following structure:

\begin{cmd}
auri@AURIMRV ~
$ ls -l example/*
-rw-r--r--    1 auri         1313 Aug  6 09:07 Dispenser.java
-rw-r--r--    1 auri         1340 Aug  6 09:07 TestDriver.java
-rw-r--r--    1 auri          923 Aug  6 09:07 VendingMachine.java
example/vending:
vending:
total 6
-rw-r--r--    1 auri         1478 Aug  6 09:49 Dispenser.class
-rw-r--r--    1 auri         1340 Aug  6 09:49 TestDriver.class
-rw-r--r--    1 auri         1253 Aug  6 09:49 VendingMachine.class
\end{cmd}

Now, from the generated \pk{.class} files, the user can create a
project using \toolname. To do this, the first step is to invoke
\toolname's graphical interface. Supposing that \toolname is
installed \rev{on} \verb+~\Tools\jabuti+, the command below causes
the invocation of its graphical interface.

\begin{cmd}
auri@AURIMRV ~/example
$ java -cp ".;..\Tools\jabuti;\
>..\Tools\jabuti\lib\BCEL.jar;\
>..\Tools\jabuti\lib\jviewsall.jar;\
>..\Tools\jabuti\lib\crimson.jar;\
>..\Tools\jabuti\lib\junit.jar" gui.JabutiGUI
\end{cmd}

Observe that the tool requires that some third-party libraries
(BCEL~\cite{Dahm01BCEB} to manipulate the Java bytecode files,
ILOG JViews~\cite{Ilog03JVIE} to visualize the \DUG, Crimson to
manipulate XML files\footnote{Such a package is required when a
java compiler under version 1.4.1 is used (considering the Sun
compiler)).}, and JUnit~\cite{JUnit02UDCA} to import test sets.)
to be included in the \pk{CLASSPATH} to allow its execution. The
current directory, (``.'' in our example), and the base directory
of the tool (\verb+..\Tools\jabuti+) are also included to allow
the correct execution of our example and the tool itself. If
desired, the user can set the \pk{CLASSPATH} environment variable
as a system variable that is initialized during the boot process.
In this case, to call \toolname's GUI, the parameter \pk{-cp} can
be omitted. Figure~\ref{fig:jabuti-initial} illustrates the
\toolname's initial window.

\begin{figure}[!ht]
\begin{center}
\includegraphics[width=0.70\textwidth]{fig/initial-screen.eps}
\caption{\label{fig:jabuti-initial}\toolname main window.}
\end{center}
\end{figure}


To create a new project, the first step is to select a base
\pk{.class} file from \pk{File $\rightarrow$ Open Class} menu. A
dialog window, as illustrated in
Figure~\ref{fig:project-select-dir} appears. From this dialog
window, the tester selects the directory where the base class file
is located and then select the base class file itself
(Figure~\ref{fig:project-select-class}). Once the base class file
is selected the tool automatically identifies the package that it
belongs (if any) and fills out the \pk{Package} field with the
package's name. The \pk{Classpath} should contains only the path
necessary to run the selected base class. In our example, since
the tool is being called inside the \verb+~\example+ directory, to
run \pk{vending.TestDriver} it is necessary that the
\pk{Classpath} field contains only the current directory, as shown
in Figure~\ref{fig:project-select-class}.

\begin{figure}[!ht]
\begin{center}
\subfigure[]{\label{fig:project-select-dir}\includegraphics[width=0.45\textwidth]{fig/project-selected-dir.eps}}\qquad
\subfigure[]{\label{fig:project-select-class}\includegraphics[width=0.45\textwidth]{fig/project-selected-class.eps}}
\caption{Open Class: (a) Selecting the directory and (b) Selecting
the main class file.}\label{fig-jabuti-open-class}
\end{center}
\end{figure}


By clicking the \pk{Open} button
(Figure~\ref{fig:project-select-class}), the \pk{Project Manager}
window, as illustrated in Figure~\ref{fig:project-manager}, will
be displayed. From the selected base class file (\pk{TestDriver}
in our example) the tool identifies the complete set of system and
non-system class files necessary to execute \pk{TestDriver}.
Currently, \toolname does not allows the instrumentation of system
class files. Therefore, the complete set of non-system class files
(user's classes) relate to \pk{TestDriver} can be instrumented.
From the \pk{Project Manager} window (left side) the user can
select the class files that will be tested. At least one class
file must be selected. In our example, we select
\pk{VendingMachine} and \pk{Dispenser} to be instrumented.
\pk{TestDriver} was not selected since it is only the driver that
allows the tester to observe whether the behavior of
\pk{VendingMachine} and \pk{Dispenser} are correct with respect to
their specification.

\begin{figure}[!ht]
\begin{center}
\includegraphics[width=0.90\textwidth]{fig/project-manager.eps}
\caption{\label{fig:project-manager}Selecting the class files to
be tested.}
\end{center}
\end{figure}


Moreover, the tester must give a name to the project being created
(\pk{vending\prjext} in our example) by clicking on the
\pk{Select} button. By clicking on the \pk{Ok} button, \toolname
creates a new project (\pk{vending.jbt} in our example),
constructs the \DUG for each method of each class under testing,
derives the complete set of testing requirements for each
criterion, calculates the weight of each testing requirement, and
presents the bytecode of a given class under testing. By creating
a project, the tester does not need to go through this entire
process again if he/she intends to test the same set of class
files. Next time, he/she can just go to \pk{File $\rightarrow$
Open Project} menu option and reload a previously saved project.

Once the project is created, the tester just need to open it and
the tool automatically detects the classes that compose it.
\toolname can be used for analyzing the coverage of a given class
file, for debugging the class file using the slice tool and for
collecting static metrics information. \toolname also allows the
visualization of the Def-Use Graph (\DUG) of each method in a
given class file as well as the visualization of the source code,
when available. Different kinds of testing reports can also be
generated considering different levels of abstraction. The
following sections describe these features in detail.
